\section{The planar equivalent-layer integral}
\label{sec:planar_eqlayer}

\subsection{Integral in the space domain}
\label{subsec:space_eqlayer_integral}

The linear system defined by equation \ref{eq:linear-system-eqlayer} represents the
discrete form of the \textit{planar equivalent layer integral}:
\begin{equation}
	d^{o}(\mathbf{r}) \approx
	d(\mathbf{r})
	= \int\limits_{-\infty}^{\infty}\int\limits_{-\infty}^{\infty}
	p(\mathbf{r}'') \, a(\mathbf{r}, \mathbf{r}'') \:\: d\mathcal{S}''
	\: , \quad z < z''_{0} \: ,
	\label{eq:eqlayer-integral}
\end{equation}
where $d\mathcal{S}'' = dx'' \, dy''$ defines an element of area; 
$\mathbf{r} = (x, y, z)^{\top}$ defines a generic observation point and
$\mathbf{r}'' = (x'', y'', z''_{0})^{\top}$ a point located on the generic horizontal 
plane at $z''_{0}$; 
$d^{o}(\mathbf{r})$ represents the observed potential field;
$d(\mathbf{r})$
is a harmonic function defining the potential field at $\mathbf{r}$ due to 
an equivalent layer with continuous physical property distribution $p(\mathbf{r}'')$; and
$a(\mathbf{r}, \mathbf{r}'')$ is a harmonic function defining the kernel.
Similarly, equation \ref{eq:transformed-field-space} assumes the form:
\begin{equation}
	f(\mathbf{r})
	= \int\limits_{-\infty}^{\infty}\int\limits_{-\infty}^{\infty}
	\tilde{p}(\mathbf{r}'') \, t(\mathbf{r}, \mathbf{r}'') \:\: d\mathcal{S}''
	\: , \quad z < z''_{0} \: ,
	\label{eq:transformed-field-integral-space}
\end{equation}
where $\tilde{p}(\mathbf{r}'')$ is a solution of the planar equivalent-layer integral (equation \ref{eq:eqlayer-integral}),
$f(\mathbf{r})$ is the transformed potential field and
$t(\mathbf{r}, \mathbf{r}'')$ defines the transformation kernel.
We stress that the observed potential field $d^{o}(\mathbf{r})$ (equation \ref{eq:eqlayer-integral})
and the transformed field $f(\mathbf{r})$ (equation \ref{eq:transformed-field-integral-space}) do not
have to be defined at the same observation point $\mathbf{r}$.

%Note that the equivalent-layer technique presumes that the potential-field 
%data can be satisfactorily described by a harmonic function.
%Besides, it also presumes that, given a potential-field data set, there is an
%equivalent layer that reproduces the observed potential-field data and can be 
%used to compute field transformations.
%It is important to stress that most of the proposed equivalent-layer methods
%considers the equivalent layer as a purely mathematical tool for processing 
%potential-field data 
%\cite[e.g.,][]{bhattacharyya_chan1977, hansen_miyazaki1984, pedersen1989, cordell1992}.
%They generally neither deduce the existence of planar equivalent layers from potential theory,
%nor discuss the underlying theoretical assumptions supporting the use of 
%a specific harmonic function $a_{ij}$ (equation \ref{eq:predicted-data-i}).
%Most of them considers that the associated physical property distribution $p(\mathbf{r}'')$ does not have
%any relationship with the true potential-field sources and also assume that
%the layer is located below the observation points and does not cross the sources;
%a condition that cannot be guaranteed in practical situations.
%What follows aims at defining planar equivalent layers as closed-form solutions 
%$p(\mathbf{r}'')$ of the planar equivalent-layer integral 
%(equation \ref{eq:eqlayer-integral}) for different kernels $a(\mathbf{r}, \mathbf{r}'')$
%and potential-field data $d^{o}(\mathbf{r})$, even for the case in which the layer crosses the sources.

\subsection{Integral in the Fourier domain}
\label{subsec:Fourier_eqlayer_integral}

Let us apply the Fourier-convolution theorem to the planar equivalent-layer integral
(equation \ref{eq:eqlayer-integral}) on a horizontal plane 
$z^{\dagger}_{0} < z^{t}_{0}$ above the sources to obtain
\begin{equation}
	\mathcal{F} [ d^{o}(\mathbf{r}^{\dagger}) ] \approx 
	\mathcal{F} [ d(\mathbf{r}^{\dagger}) ] =
	\mathcal{F} [ p(\mathbf{r}'') ] \, 
	\mathcal{F} [ a(\mathbf{r}^{\dagger}, \mathbf{r}'') ] \: ,
	\label{eq:eqlayer-integral-Fourier}
\end{equation}
where
\begin{equation}
	\mathcal{F} [ a(\mathbf{r}^{\dagger}, \mathbf{r}'') ] =
	\alpha_{a} \, 
	\mathcal{F} [ \Phi(\mathbf{r}^{\dagger}, \mathbf{r}'') ] \: ,
	\quad \left| \alpha_{a} \right| \ne 0 \: ,
	\label{eq:kernel-Fourier}
\end{equation}
with $\alpha_{a}$ being a generic complex or real function of the wavenumbers $k_{x}$ and $k_{y}$.
This function may be defined in terms of the Fourier-domain differential operator 
$\alpha(\ell,\hat{\mathbf{u}})$ (equation \ref{eq:diff-operators-generic}), using 
$\ell$ and $\hat{\mathbf{u}}$ not necessarily equal to
that used to define $\alpha_{d}$ (equation \ref{eq:observed-data-Fourier}).
Similarly, we get by applying the Fourier-convolution theorem to equation \ref{eq:transformed-field-integral-space}
that
\begin{equation}
	\mathcal{F} [ f(\mathbf{r}^{\ddagger}) ] =
	\mathcal{F} [ \tilde{p}(\mathbf{r}'') ] \, 
	\mathcal{F} [ t(\mathbf{r}^{\ddagger}, \mathbf{r}'') ] \: ,
	\quad z^{\ddagger}_{0} < z''_{0} \: ,
	\label{eq:transformed-field-integral-Fourier}
\end{equation}
where $\mathbf{r}^{\ddagger} = (x^{\ddagger}, y^{\ddagger}, z^{\ddagger}_{0})^{\top}$ is a generic point not
necessarily equal to $\mathbf{r}^{\dagger}$, 
$\mathcal{F} [ \tilde{p}(\mathbf{r}'') ]$ is the Fourier transform of a particular solution $\tilde{p}(\mathbf{r}'')$ 
of the planar equivalent-layer integral (equation \ref{eq:eqlayer-integral}) and
\begin{equation}
	\mathcal{F} [ t(\mathbf{r}^{\ddagger}, \mathbf{r}'') ] =
	\alpha_{t} \, 
	\mathcal{F} [ \Psi(\mathbf{r}^{\ddagger}, \mathbf{r}'') ] \: ,
	\quad \left| \alpha_{t} \right| \ne 0 \: ,
	\label{eq:kernel-transformation-Fourier}
\end{equation}
defines the kernel of the potential-field transformation in the Fourier domain,
with $\alpha_{t}$ being another generic complex or real function of the wavenumbers $k_{x}$ and $k_{y}$.

\subsection{Auxiliary functions}

The subsequent sections aims at presenting solutions for the planar equivalent-layer integral (equation \ref{eq:eqlayer-integral}).
To do that, we need first consider a function $\tau(\mathbf{r}'')$ defined on a horizontal 
plane $z''_{0}$ that may be located above (i.e., $z''_{0} < z^{t}_{0}$) or cross 
the sources (i.e., $z''_{0} > z^{t}_{0}$), with 2DFT given by
\begin{equation}
	\mathcal{F} \left[ \tau(\mathbf{r}'') \right] = 
	\begin{cases}
		\mathcal{F} \left[ U(\mathbf{r}^{\dagger}) \right] 
		e^{\, k \, (z''_{0} - z^{\dagger}_{0})} & , 
		\quad 0 < k \le k_{\varepsilon} \\
		0 & , \quad k > k_{\varepsilon} 
	\end{cases} \: ,
	\label{eq:tau-Fourier}
\end{equation}
which is defined in terms of a fixed and possibly very large truncating wavenumber $k_{\varepsilon}$
and leads to the auxiliary potential
\begin{equation}
	V(\mathbf{r}, z''_{0}) = 
	\int\limits_{-\infty}^{\infty} \int\limits_{-\infty}^{\infty}
	\tau(\mathbf{r}'') \:
	\Phi(\mathbf{r}, \mathbf{r}'') 
	\:\: d\mathcal{S}'' \: , 
	\quad z < z^{t}_{0} \: , \quad z < z''_{0} \: .
	\label{eq:potential-V-space}
\end{equation}

Note that the function $\tau(\mathbf{r}'')$ obtained by computing
the inverse Fourier transform of 
$\mathcal{F} \left[ \tau(\mathbf{r}'') \right]$ (equation \ref{eq:tau-Fourier})
approximates the potential field obtained by 
computing the downward continuation of $U(\mathbf{r}'')$ (equation \ref{eq:potential-U-space}) 
to points that may be located inside the sources. 
It is well-known, however, that downward continuation can only be computed at 
source-free points, where the Laplace's equation applies. This means that the
continuation plane is limited by the top $z^{t}_{0}$ of the sources, otherwise the continued field 
presents an oscillatory character \citep[e.g.,][]{roy1967, fedi_florio2002} that may
prevent $\mathcal{F} \left[ \tau(\mathbf{r}'') \right]$ to be square integrable.
The truncating wavenumber $k_{\varepsilon}$ (equation \ref{eq:tau-Fourier}) 
overcomes this problem, as pointed out by \citet{bhattacharyya1967},
and guarantees that the Fourier-convolution theorem 
can be applied to the auxiliary potential $V(\mathbf{r}, z''_{0}, k_{\varepsilon})$
(equation \ref{eq:potential-V-space}) on a horizontal plane $z^{\dagger}_{0} < z^{t}_{0}$ 
above the sources to obtain
\begin{equation}
	\mathcal{F} [ V(\mathbf{r}^{\dagger}, z''_{0}) ]
	= \tfrac{k}{\alpha_{\Phi}} \, \mathcal{F} [ \tau(\mathbf{r}'') ] \:
	\mathcal{F} [ \Phi(\mathbf{r}^{\dagger}, \mathbf{r}'') ] \: ,
	\label{eq:potential-V-Fourier}
\end{equation}
where $\mathcal{F} [ \Phi(\mathbf{r}^{\dagger}, \mathbf{r}'') ]$ is given by
equation \ref{eq:Phi-Fourier}. 

We get by substituting 
$\mathcal{F} \left[ \Phi(\mathbf{r}^{\dagger}, \mathbf{r}'') \right]$
(equation \ref{eq:Phi-Fourier}) and 
$\mathcal{F} \left[ \tau(\mathbf{r}'') \right]$
(equation \ref{eq:tau-Fourier}) into 
$\mathcal{F} [ V(\mathbf{r}^{\dagger}, z''_{0}) ]$ (equation \ref{eq:potential-V-Fourier}) that
\begin{equation}
	\mathcal{F} [ V(\mathbf{r}^{\dagger}, z''_{0}) ] = 
	\begin{cases}
		\mathcal{F} [ U(\mathbf{r}^{\dagger}) ] &
		, \quad 0 < k \le k_{\varepsilon} \\
		0 & , \quad k > k_{\varepsilon} 
	\end{cases} \: ,
	\label{eq:potential-V-Fourier-convolution}
\end{equation}
which shows that the Fourier transforms of 
$U(\mathbf{r}^{\dagger})$ (equation \ref{eq:potential-U-Fourier}) and 
$V(\mathbf{r}^{\dagger}, z''_{0})$ (equation \ref{eq:potential-V-space}) 
on the horizontal plane $z^{\dagger}_{0}$ located above the 3D sources coincide up 
to the truncating wavenumber $k_{\varepsilon}$.

At last, it can be easily verified that, if we let the truncating wavenumber $k_{\varepsilon}$ 
goes to infinite for the case in which $z''_{0} < z^{t}_{0}$, we get 
\begin{equation}
	\lim_{k_{\varepsilon}\to\infty} \mathcal{F} \left[ \tau(\mathbf{r}'') \right] = 
	\mathcal{F} \left[ U(\mathbf{r}'') \right] \: , \quad z''_{0} < z^{t}_{0} \: ,
	\label{eq:lim-kepsilon-tau-Fourier}
\end{equation}
\begin{equation}
	\lim_{k_{\varepsilon}\to\infty} \mathcal{F} \left[ V(\mathbf{r}^{\dagger}, z''_{0}) \right] = 
	\mathcal{F} \left[ U(\mathbf{r}^{\dagger}) \right] \: , \quad z^{\dagger}_{0} < z''_{0} < z^{t}_{0} \: ,
	\label{eq:lim-kepsilon-potential-V-Fourier}
\end{equation}
and, consequently, that
\begin{equation}
	\lim_{k_{\varepsilon}\to\infty} V(\mathbf{r}, z''_{0}) = 
	U(\mathbf{r}) \: , \quad z < z''_{0} < z^{t}_{0} \: ,
	\label{eq:lim-kepsilon-potential-V-space}
\end{equation}
which is in perfect agreement with the solutions of the Dirichlet's and Neumann's
problems on a plane (equations \ref{eq:U-Dirichlet}, \ref{eq:U-Neumann} and \ref{eq:U-Dirichlet-Neumann-Fourier}).